\notechapter{N-20230313}{Category Theory Continued: Posets, Groups as One-Object Categories, Products, and Monomorphisms}

\section{What this note continues}

Elementary category theory changes the intuition of what an arrow may be.  A morphism is often introduced as a structure-preserving function, but arrows may instead be order relations, group elements, inclusions, or other constructions determined by the category.

The main examples are:
\[
\begin{gathered}
  \text{posets as categories},\quad
  \text{groups as one-object categories},\\
  \text{products of categories},\quad
  \text{monomorphisms and epimorphisms}.
\end{gathered}
\]

\section{Posets are categories}

Let $(P,\le)$ be a partially ordered set.  We can build a category, still denoted $P$, as follows.  The objects are the elements of $P$.  For two objects $p,q\in P$, define
\[
  \operatorname{Hom}(p,q)=
  \begin{cases}
  \{*\},&p\le q,\\
  \varnothing,&p\not\le q.
  \end{cases}
\]
Thus there is at most one arrow from $p$ to $q$.  The identity arrow exists because $p\le p$.  Composition exists because order is transitive: if $p\le q$ and $q\le r$, then $p\le r$.

This example is surprisingly important.  It shows that arrows in a category need not be functions.  In the poset category, the arrow $p\to q$ is the fact that $p\le q$.

\section{A small poset example}

Suppose $P$ has four objects $a,b,c,d$ with relations
\[
  a\le b,\qquad a\le c\le d,\qquad a\le d,
\]
and identities at every object.  The corresponding category has arrows exactly for these order relations.  In a Hasse diagram, upward paths encode arrows, but the category contains the composite arrow as well.  For example, from $a\le c$ and $c\le d$, the composite gives the arrow $a\to d$.

The same symbol $\le$ can be treated as an arrow.  This is not a metaphor; it satisfies the category axioms.

\section{Isomorphisms in a poset category}

An isomorphism in a category is an arrow
\[
  f:A\to B
\]
for which there exists
\[
  g:B\to A
\]
such that
\[
  g\circ f=\operatorname{id}_A,\qquad
  f\circ g=\operatorname{id}_B.
\]
In a poset category, an arrow $p\to q$ means $p\le q$.  An isomorphism between $p$ and $q$ would require arrows both ways:
\[
  p\le q,\qquad q\le p.
\]
By antisymmetry of a partial order, this forces $p=q$.  Therefore no two distinct objects of a poset category are isomorphic.

The result becomes obvious only after translating categorical isomorphism back into the order axioms.

\section{Groups as one-object categories}

A group $G$ can be regarded as a category $\mathcal G$ with one object $*$.  The arrows from $*$ to itself are the elements of $G$:
\[
  \operatorname{Hom}_{\mathcal G}(*,*)=G.
\]
Composition of arrows is group multiplication.  The identity arrow is the identity element $e\in G$.  Inverses of arrows are group inverses.

Thus a group is exactly a category with one object in which every arrow is an isomorphism.  This is not a trick.  It says that a group can be understood as the collection of all symmetries of one object, with composition as the operation.

\section{Why this example matters}

If one first learns category theory, a group may look like a special category: one object, only reversible arrows.  This reverses the usual learning order.  The note quotes the idea that a group is what one gets by collecting all symmetries of an object.  From the categorical viewpoint, a group is the extreme case where there is only one object and every transformation is invertible.

The slogan is:
\[
  \text{monoid} = \text{one-object category},\qquad
  \text{group} = \text{one-object groupoid}.
\]
A groupoid is a category in which every arrow is invertible.

\section{Products of categories}

Given categories $\mathcal C$ and $\mathcal D$, their product category $\mathcal C\times\mathcal D$ has objects
\[
  (C,D),\qquad C\in\mathcal C,\ D\in\mathcal D,
\]
and morphisms
\[
  (f,g):(C,D)\to(C',D')
\]
where $f:C\to C'$ is a morphism in $\mathcal C$ and $g:D\to D'$ is a morphism in $\mathcal D$.  Composition is componentwise:
\[
  (f',g')\circ(f,g)=(f'\circ f,\ g'\circ g).
\]
The identity on $(C,D)$ is $(\operatorname{id}_C,\operatorname{id}_D)$.

This is the categorical version of taking a product structure coordinate by coordinate.

\section{The opposite category}

For any category $\mathcal C$, the opposite category $\mathcal C^{op}$ has the same objects as $\mathcal C$, but every arrow is reversed:
\[
  f:A\to B\text{ in }\mathcal C
  \quad\longleftrightarrow\quad
  f^{op}:B\to A\text{ in }\mathcal C^{op}.
\]
Composition is reversed accordingly.  The opposite category is a formal way to turn statements about maps out of an object into statements about maps into an object.

For posets, passing to the opposite category corresponds to reversing the order.  If $P$ is viewed as a category, then $P^{op}$ is the poset with
\[
  p\le_{op}q\quad\Longleftrightarrow\quad q\le p.
\]

\section{Monomorphisms and epimorphisms}

In ordinary set theory, an injective function is one that does not identify distinct elements.  Category theory cannot always talk about elements, so it defines a left-cancellation property instead.

A morphism $f:X\to Y$ is a monomorphism if for every pair of arrows $g,h:Z\to X$,
\[
  f\circ g=f\circ h\quad\Longrightarrow\quad g=h.
\]
It is an epimorphism if for every pair of arrows $g,h:Y\to Z$,
\[
  g\circ f=h\circ f\quad\Longrightarrow\quad g=h.
\]
In $\mathbf{Set}$, monomorphisms are exactly injective functions, and epimorphisms are exactly surjective functions.  In other categories, this may fail, which is why the categorical definitions are more fundamental.

\section{The lesson of the note}

The guiding idea is to stop identifying arrows with functions.  In a poset they are inequalities, in a group-as-category they are group elements, in a product category they are pairs, and in an opposite category they are formally reversed.  Monomorphism and epimorphism then express injective-like and surjective-like behavior without mentioning elements.

The transition is from element-based thinking to arrow-based thinking.  This is the first real shift in category theory.
